% compendium/execution-market-making/order-books.tex
% Entries migrated on 2026-10-16 from the legacy order-lifecycle, limit-order-book-mechanics and queue-dynamics files and
% the order-flow-imbalance file (slice 33). Owners: the section overview execution-market-making (lifecycle, the book,
% limit orders as options), market-making-lob (queues, queue imbalance) and lob-deep-learning (order-flow imbalance).
% Dispositions in notes/migration-slices.md; the legacy files verbatim in notes/artefacts/<owner>/.

\section{Orders, limit order books and queues}

\begin{qfmechanism}[Order lifecycle and messages: order states, insert, modify, cancel and execute, private and public channels, and the instructions that change priority]\label{mech:order-lifecycle}
\qffield{Statement} An order book is empty at the start of the day; orders are stored while they do not match a resting order, and two orders match when their prices are compatible, each match producing one public trade message and two private messages, one to the owner of each order \citep[chap.~3, sec.~3.3.1, p.~217]{alpha-laruelle}. Members usually reach the matching engine through two channels: a private one, on which they send their orders and receive the messages that concern them, and a public one, on which everyone connected sees the transactions and the state of the book; the private channel is generally the faster (p.~218). At any time a trader can insert, modify or cancel an order and attach tags to it, such as a time in force (immediate or cancel, good till day, good till time) or a displayed quantity, which creates an iceberg order (p.~218). An order then leaves the book by execution against an incoming order that crosses it, by cancellation, or by expiry; an iceberg's hidden part usually has lower time priority than displayed orders at the same price submitted later, and an immediate-or-cancel order that matches nothing leaves no visible trace \citep[chap.~3, sec.~3.2.2, p.~55]{catalogue-0528}. The states an order passes through are those a broker or venue reports for it; in the FIX~4.4 message standard they include pending new, new, partially filled, filled, done for day, cancelled, replaced, pending cancel, pending replace, stopped, rejected, suspended and expired \citep[OrdStatus, tag~39]{onixs-fix44-ordstatus}.

\qffield{Notation} \emph{Insert} adds displayed or hidden quantity at a price; \emph{modify} changes price, quantity or instructions; \emph{cancel} removes resting quantity; \emph{execute} is a match against an incoming marketable order. \emph{Working} (resting) and \emph{part-filled} orders have quantity still governed by their instructions; filled, cancelled, expired and rejected orders are terminal.

\qffield{Conditions} Priority after a modification is a venue rule, not a general one. On CME Globex under FIFO matching an order loses its priority and is requeued if its quantity is increased, its price changed or its account changed \citep[page ``Supported Matching Algorithms'', updated 11 February 2025]{ex-cme-matching}; a reduction in quantity is not on that list. Under pro-rata matching, which Bouchaud et al.\ describe as common in short-rate futures, time does not break ties at a price, and each resting order receives a share of an incoming order in proportion to its size \citep[chap.~3, sec.~3.2.1, p.~54]{catalogue-0528}. A rule that applies to one venue, order type or matching algorithm must be stated with that venue and its rule date.

\qffield{Practice} Limit and stop instructions let a trader who is absent take part: many traders use limit orders to tell their brokers the conditions under which they would trade, and brokers then watch the market on their behalf \citep[chap.~4, sec.~4.4.5, pp.~77--78]{alpha-harris}. A stop order becomes an ordinary order once triggered and ``remains valid even if price crosses back over the stop price'', so it can fill well beyond its stop price in a fast market (Harris's cotton example: a stop at 70 cents sold at 66.90; sec.~4.5, p.~78). For a replay, an execution record therefore needs the acknowledgement, every state change and the channel timestamps of each order, not only its fills; the public feed alone cannot show which resting order was the replayer's own.

\qffield{Pitfalls} The legacy notes said a modification ``loses priority when price worsens or size increases''. On CME's FIFO algorithm any price change loses priority, an improvement included; the legacy rule matches no venue rule read here. Their five ``order states'' (pending, accepted, working, part-filled, terminal) are a simplification of the reported states above, and they named minimum-quantity and cancel-on-disconnect instructions, which were not checked against any venue's rules and are left as a backlog item.

\qffield{See also} \hyperref[def:limit-order-book]{limit order book}; \hyperref[conv:cme-globex-trading]{CME Globex trading conventions} (matching algorithms); \hyperref[def:order-exposure-decisions]{order exposure decisions} (hidden and iceberg orders); \hyperref[meth:automated-trading-operational-controls]{operational controls} (unknown order state); \hyperref[data:market-data-levels]{market data levels} (Level~III order messages).

\qffield{Further reading} The pathway's required readings on replaying order books (Huang, Lehalle and Rosenbaum; Noble, Rosenbaum and Souilmi) are not read here. The legacy file is kept verbatim in \nolinkurl{notes/artefacts/execution-market-making/}.

\qffield{Index} order lifecycle; order state; OrdStatus; pending new; partially filled; insert; modify; cancel; execute; private channel; public feed; time in force; immediate or cancel; iceberg order; loss of priority; requeue; stop order; absent traders

\qffield{Status} Verification: partly checked 2026-10-16 (against the text layers of local copies of \texttt{alpha-laruelle}, 2018 World Scientific printing, pp.~217--218; \texttt{catalogue-0528}, Cambridge 2018, pp.~54--55; \texttt{alpha-harris}, 2002 Oxford printing, pp.~77--78; the CME page as cited in slice 19, not reread; and the OnixS rendering of FIX~4.4, read 2026-10-16). Stability: draft. Currency: FIX~4.4 states and the CME rule as of 2025; other venues not checked.
\end{qfmechanism}

\begin{qfdefinition}[Limit order book: best quotes, limit-price placement, visibility, the three event types and how prices change, and quote-driven requests for quotes]\label{def:limit-order-book}
\qffield{Statement} An order that matches on arrival is a market order and one that does not is a limit order; a \emph{limit order book} is the set of all limit orders for an asset on one platform at a given time, its best buy price is the bid $b(t)$, its best sell price the ask $a(t)>b(t)$, with mid-price $m(t)=\tfrac12[a(t)+b(t)]$ and spread $s(t)=a(t)-b(t)$, and the book is a set of queues, one per price level \citep[chap.~3, secs.~3.1.3--3.1.4, pp.~45--47]{catalogue-0528}. Limit orders supply liquidity and market orders demand it: patient traders meet impatient ones (p.~43). Harris classifies a limit order by its price relative to the market \citep[chap.~4, sec.~4.4.1, Table~4-1, pp.~73--74]{alpha-harris}: \emph{marketable} (a buy at or above the best offer), \emph{making a new market} (between the best bid and offer), \emph{at the market} (at the best bid or offer), and \emph{behind} or, further away, \emph{far from} the market; aggressiveness falls in that order. Open-book markets display their books to all traders, closed-book markets do not, and some let traders hide their orders \citep[chap.~5, p.~109]{alpha-harris}.

\qffield{Conditions} Book dynamics are driven by three kinds of event: insertions add liquidity, while cancellations and transactions (by market or marketable limit orders) remove it, and trades can occur only at the best limits \citep[chap.~3, sec.~3.3.1, p.~219]{alpha-laruelle}. The bid rises when a buy limit order arrives inside the spread and falls when sells exhaust the volume at the bid; a sell market order larger than the bid queue walks down the book and trades at several prices \citep[chap.~3, sec.~3.1.8, p.~51; Table~3.1, p.~53]{catalogue-0528}. Request-for-quote systems differ: a trader asks several market makers, each free to answer, and modern systems aggregate the answers into a book-like view; Lehalle and Laruelle note fixed-income markets' move towards all-to-all RFQs that look like central limit order books \citep[chap.~3, p.~216]{alpha-laruelle}. Interpretation: an aggregated RFQ screen shows conditional, request-specific prices, not firm resting orders, so its depth is not comparable with a book's.

\qffield{Practice} Bouchaud et al.\ measure prices relative to the same-side best quote, $d=b(t)-p_x$ for a buy and $p_x-a(t)$ for a sell, to aggregate order flow across traders (sec.~3.1.6, p.~50). The spread and mid-price are measured as in \hyperref[def:spread-measures]{spread measures}, which this entry does not repeat.

\qffield{Pitfalls} The legacy taxonomy's ``far from the market: low immediacy cost'' reverses Harris's point: such orders give a better price if they fill but rarely fill. Its ``top-of-book'' and ``depth book'' visibility levels are data products, defined in \hyperref[data:market-data-levels]{market data levels}. A displayed book is not all liquidity: in Bouchaud et al.'s 2015 NASDAQ sample, 3 to 5 percent of market orders matched a hidden order inside the spread (Table~4.1, p.~60).

\qffield{See also} \hyperref[mech:order-lifecycle]{order lifecycle}; \hyperref[res:limit-order-options]{limit orders as trading options}; \hyperref[mod:queue-reactive-intensities]{queue-reactive intensities}; \hyperref[def:spread-measures]{spread measures}; \hyperref[def:tick-size]{tick size}; \hyperref[def:floor-automated-trading-systems]{floor and automated trading systems}; \hyperref[def:fixed-income-market-structure]{fixed-income market structure} (RFQ protocols).

\qffield{Further reading} The legacy file is kept verbatim in \nolinkurl{notes/artefacts/execution-market-making/}.

\qffield{Index} limit order book; market order; limit order; bid; ask; mid-price; marketable limit order; at the market; behind the market; making a new market; open book; closed book; insertion; cancellation; transaction; walking the book; quote-relative price; request for quote; RFQ; central limit order book

\qffield{Status} Verification: partly checked 2026-10-16 (against the text layers of local copies of \texttt{catalogue-0528}, Cambridge 2018, pp.~43--53 and~60; \texttt{alpha-harris}, 2002 Oxford printing, pp.~73--74 and~109; \texttt{alpha-laruelle}, 2018 World Scientific printing, pp.~216--219). Stability: draft. Currency: definitions timeless; the hidden-order shares are NASDAQ 2015.
\end{qfdefinition}

\begin{qfresult}[Standing limit orders are trading options: what sets their value, the two risks of using them, and why fills are adversely selected]\label{res:limit-order-options}
\qffield{Statement} Standing limit orders are options to trade: a sell limit order is a call that lets others buy when they want to, a buy limit order a put that lets them sell, and the strike is the limit price \citep[chap.~4, sec.~4.4.2, p.~75]{alpha-harris}. They are options but not option contracts: their writers grant them free to the market and can cancel them, and anyone can exercise one with a market or marketable limit order (p.~75). Their value to others rises with how close the limit price is to the market, how long the order will stand, and volatility; to limit that value, traders place orders further from the market when it is volatile, so spreads are wide in volatile markets (pp.~75--76). Bouchaud et al.\ make the same point, adding that limit orders ``typically earn profits when prices mean-revert but suffer losses when prices trend'' \citep[chap.~3, sec.~3.1.3, p.~46]{catalogue-0528}. The legacy formalisation, derived here: for an order of size $q$ at limit $L$, a contra trader with value $V_t$ exercises only when that is worth while, so the exercise value is $q(V_t-L)^+$ for a sell order and $q(L-V_t)^+$ for a buy order, while the order is live.

\qffield{Conditions} The option view needs a contra trader whose value differs from the limit price and who can act before the writer cancels; the payoff is the exerciser's surplus, not the writer's loss on a market that later moves. Two risks follow \citep[chap.~4, sec.~4.4.4, p.~77]{alpha-harris}: \emph{execution uncertainty}, when prices move away and the order fails to trade (Harris's example: a buyer who chased the price paid 9 cents more than a market order would have, p.~77), and \emph{ex post regret}, when prices move through the limit and the order fills just before the market moves against the new position. Adverse selection, trading with better-informed traders, is ``the most important cause of ex post regret''. Placing orders further away lowers regret and the fill probability; placing them closer does the reverse.

\qffield{Practice} Interpretation, from the passages above: fills are not a random sample of the times an order was available, because informed or faster traders exercise it when it has become cheap, so an evaluation of a passive strategy must measure the price move after its own fills (as in the realised spread of \hyperref[def:spread-measures]{spread measures}), not only its fill rate. In a continuous order-driven market the equilibrium spread makes the marginal trader indifferent between limit and market orders; that result, with its conditions, is \hyperref[res:spread-determinants]{spread determinants}, which this entry does not restate.

\qffield{Pitfalls} The legacy proposition gave the payoff ``to the market'' without saying it is exercised at the other trader's discretion and that the writer may cancel; its ``timing option'' mechanism added cancellation speed and monitoring speed, which match Harris's argument about volatility and time to cancel (\hyperref[res:spread-determinants]{spread determinants}) but were not stated as such in the pages read here.

\qffield{See also} \hyperref[def:limit-order-book]{limit order book}; \hyperref[res:spread-determinants]{spread determinants}; \hyperref[mod:static-limit-order-depth]{static limit-order depth}; \hyperref[res:latency-arbitrage-arms-race]{latency arbitrage}; \hyperref[mod:glosten-milgrom]{Glosten--Milgrom}.

\qffield{Further reading} The legacy blocks are kept verbatim in \nolinkurl{notes/artefacts/execution-market-making/} (limit-order-book file) and \nolinkurl{notes/artefacts/market-making-lob/} (queue file).

\qffield{Index} limit order as option; free trading option; call option; put option; option value of a limit order; execution uncertainty; ex post regret; adverse selection; chasing the price; equilibrium spread; timing option

\qffield{Status} Verification: partly checked 2026-10-16 (against the text layers of local copies of \texttt{alpha-harris}, 2002 Oxford printing, pp.~75--77, and \texttt{catalogue-0528}, Cambridge 2018, p.~46; the payoff expression derived here). Stability: draft. Currency: timeless.
\end{qfresult}

\begin{qfdefinition}[Queue imbalance at the best quotes: two normalisations, its predictive power for the next price move, and its limits]\label{def:queue-imbalance}
\qffield{Statement} With $V_b$ and $V_a$ the volumes at the best bid and ask, Bouchaud et al.\ define the \emph{queue imbalance}
\[ I=\frac{V_b}{V_b+V_a}\in(0,1),\qquad\text{equivalently}\qquad QI=\frac{V_b-V_a}{V_b+V_a}=2I-1\in(-1,1), \]
invariant to rescaling both volumes; $I$ near 1 (a much larger bid queue) suggests net buying pressure and an upward move \citep[chap.~7, sec.~7.2, eq.~(7.1), p.~119]{catalogue-0528}. For a large-tick stock whose spread is one tick, the price changes only when one best queue empties, so the next move is a race between the two queues (sec.~7.1, p.~118). Estimating from random times the share of races that end with an up move, they find ``a strong, monotonic correlation between the queue imbalance and the direction of the next price move'' (p.~120). Lehalle and Laruelle state the same fact in queue-reactive terms: ``when the first queue is smaller than its opposite queue, it will most probably be consumed'' \citep[chap.~3, sec.~3.3.2, pp.~223--224]{alpha-laruelle}.

\qffield{Conditions} \emph{Sample.} Bouchaud et al.'s figure uses two large-tick NASDAQ stocks, INTC and CSCO, in 2015 (chap.~4, p.~58; Table~4.1, p.~60); Lehalle and Laruelle's statistics are for a large-tick stock with an average spread of about 1.3 ticks (p.~227). \emph{Horizon.} The prediction is for the next mid-price move, and Lehalle and Laruelle add that ``the amplitude of this future price move is far lower than a tick'' in most cases, so that no rewarding arbitrage can be built on it alone (p.~224). \emph{Large and small ticks.} The race argument needs a spread pinned at one tick; for small-tick stocks the spread can be filled inside and the argument does not apply as stated.

\qffield{Practice} Checked (\nolinkurl{order-flow-checks.py}): $QI=2I-1$ to machine precision, and both are unchanged when both volumes are scaled. Interpretation: as a model feature it must be bound to a stated forecast horizon and validated on later days and other stocks, since hidden orders, replenishment after depletion and the next queue's size all enter after the next move.

\qffield{Pitfalls} The legacy notes defined queue imbalance twice, once as $(B_t-A_t)/(B_t+A_t)$ and once in words, with a second claim (``a large opposite queue can make the next price move less likely to reverse immediately'') found in neither source; it is not adopted. Queue imbalance is a state of displayed volume, distinct from the event-based \hyperref[def:order-flow-imbalance]{order-flow imbalance} and from signed trade volume.

\qffield{See also} \hyperref[mod:queue-reactive-intensities]{queue-reactive intensities}; \hyperref[def:order-flow-imbalance]{order-flow imbalance}; \hyperref[def:tick-size]{tick size} (large and small ticks).

\qffield{Further reading} The pathway's readings on imbalance-based placement (Lehalle and Mounjid) and on why imbalance predicts prices (Pulido, Rosenbaum and Sfendourakis, an optional branch) are not read here; Bouchaud et al.'s models of the race (secs.~7.3--7.5) are not summarised. The legacy blocks are kept verbatim in \nolinkurl{notes/artefacts/market-making-lob/} and \nolinkurl{notes/artefacts/lob-deep-learning/}.

\qffield{Index} queue imbalance; volume imbalance; best-queue imbalance; race to the bottom; next price move; large-tick stock; worst-kept secret of high-frequency trading

\qffield{Status} Verification: partly checked 2026-10-16 (against the text layers of local copies of \texttt{catalogue-0528}, Cambridge 2018, pp.~58, 60 and 118--120, and \texttt{alpha-laruelle}, 2018 World Scientific printing, pp.~223--224 and~227; the identity computed). Stability: draft. Currency: evidence from NASDAQ 2015 and an unnamed European large-tick stock.
\end{qfdefinition}

\begin{qfmodel}[Queue-reactive order-book intensities: the Poisson baseline and why it fails, intensities conditioned on queue sizes, the arrival--departure ratio, and the reset that limits mean reversion]\label{mod:queue-reactive-intensities}
\qffield{Statement} Conditioning a variable means specifying the contexts in which it is ``less random''; for an order book, the probability of an insertion at the first limit when it is large, rather than of any insertion anywhere \citep[chap.~3, sec.~3.3.2, pp.~219--220]{alpha-laruelle}. The baseline models three intensities, of insertions $\lambda^+$, cancellations $\lambda^c$ and transactions $\lambda^x$, as homogeneous Poisson processes independent of the book; each is estimated by starting at many times $t$, recording the waiting time $\tau-t$ to the next event of that type and averaging, which estimates $1/\lambda$ (p.~221). Simulated with those constant intensities, the book's queue sizes are ``very far away from the reality'' (Figure~3.11, p.~222). Huang, Lehalle and Rosenbaum's \emph{Model~I} keeps one intensity per event type and per size of the queue concerned, measured in average event sizes, by filing each waiting time under the queue size at the start time; Models~II and~III also condition on the sizes of the neighbouring queues, for example $\lambda^c(4,2,1)$ for a first bid of 4 with a second bid of 2 and a best ask of 1 (pp.~222--223).

\qffield{Notation} The \emph{arrival--departure} (A/R) ratio of a queue is its insertion intensity over the sum of its cancellation and transaction intensities; Lehalle and Laruelle read a ratio above one as a queue that grows on average and below one as one that depletes (p.~224). AES: average event size, the average size of an insertion, cancellation or trade.

\qffield{Conditions} In their figures the first queue's ratio is below one at every size and the second's above one, and the book is stable only because, when the first queue empties, the second is promoted and starts to deplete (pp.~224--225; the book is ergodic, n.~13). Two regularities explain why conditioning helps: a first queue smaller than the opposite one is usually consumed, and a large first queue brings more cancellations to the second (pp.~223--224). The statistics are for a large-tick stock with a spread of about 1.3 ticks; with small ticks a queue carries less meaning (p.~227). Interpretation, checked on a birth--death toy (\nolinkurl{order-flow-checks.py}): with intensities that do not depend on size, a ratio of 1.2 gives an upward drift of $+0.09$ per event and 0.8 a drift of $-0.11$; with size-dependent intensities the ratio has to be read at each size.

\qffield{Practice} Liquidity dynamics fitted this way give simulated prices whose volatility is ``far too low'': a newly promoted first queue is large, which pulls the price back, so mean reversion is too strong (pp.~225--226). Lehalle and Laruelle suggest Hawkes processes to add memory, or a reset: with some probability $\theta$, when a first queue depletes, the book is redrawn from the database as if participants accepted the new price (p.~226). The \emph{starting time} of the waiting-time estimator matters: from every second, from just after any event, or just after a trade gives different estimates (p.~221, n.~12). Checked (\nolinkurl{order-flow-checks.py}): for a clustered (Hawkes) sample the mean wait from uniform start times is 0.594 against a mean inter-event time of 0.405; the two agree in expectation only for a Poisson process.

\qffield{Pitfalls} The queue-reactive model is a conditional Markov model with exponential waiting times; the guide's catalogue notes that it misses clustered latency races and a trader's longer-lived impact, so a simulator built on it must reproduce both book statistics and policy costs before its fills are trusted. The legacy ``latency race boundary'' (latency matters most when a trader can change queue position, less with wide ticks, deep queues or orders not competing for priority) is interpretation; Lehalle and Laruelle say a small tick shifts attention ``from queuing issues to limit price issues'' (chap.~1, p.~89), and Bouchaud et al.\ that activity at the best quotes is faster for large-tick stocks, reflecting the importance of queue position (chap.~4, p.~59).

\qffield{See also} \hyperref[def:queue-imbalance]{queue imbalance}; \hyperref[mod:hawkes-order-flow]{Hawkes processes for order flow}; \hyperref[def:order-event-point-processes]{point processes and intensities}; \hyperref[mod:static-limit-order-depth]{static limit-order depth}; \hyperref[meth:passive-order-split]{passive order split} (queue position across venues).

\qffield{Further reading} Huang, Lehalle and Rosenbaum (2015) and Cont, Stoikov and Talreja (2010), required readings of \texttt{market-making-lob}, are not read here: the account is Lehalle and Laruelle's. The legacy blocks and the order-flow checks are in \nolinkurl{notes/artefacts/market-making-lob/} and \nolinkurl{notes/artefacts/market-simulation/}.

\qffield{Index} queue-reactive model; Model I; conditioning; intensity; insertion; cancellation; transaction; homogeneous Poisson; waiting time; average event size; arrival--departure ratio; queue promotion; ergodic order book; mean reversion; book reset; latency and queue position

\qffield{Status} Verification: partly checked 2026-10-16 (against the text layer of a local copy of \texttt{alpha-laruelle}, 2018 World Scientific printing, chap.~3, pp.~219--227, and chap.~1, p.~89; \texttt{catalogue-0528}, p.~59; the drift and waiting-time illustrations computed). Stability: draft. Currency: the model as described in 2018.
\end{qfmodel}

\begin{qfdefinition}[Order-flow imbalance: Cont, Kukanov and Stoikov's event-based measure at the best quotes, against signed trade volume, and the family of imbalance features]\label{def:order-flow-imbalance}
\qffield{Statement} For consecutive best-quote states $(P^B_n,q^B_n,P^A_n,q^A_n)$, Cont, Kukanov and Stoikov define the contribution of the $n$th event to supply and demand at the best quotes,
\[ e_n=\mathbf 1\{P^B_n\ge P^B_{n-1}\}q^B_n-\mathbf 1\{P^B_n\le P^B_{n-1}\}q^B_{n-1}-\mathbf 1\{P^A_n\le P^A_{n-1}\}q^A_n+\mathbf 1\{P^A_n\ge P^A_{n-1}\}q^A_{n-1}, \]
so that at an unchanged bid $e_n$ is the change in bid size, a higher bid contributes the new bid size and a lower bid minus the old one, and the ask side enters with opposite signs; the \emph{order-flow imbalance} over $[t_{k-1},t_k]$ is $\mathrm{OFI}_k=\sum_{n=N(t_{k-1})+1}^{N(t_k)}e_n$, which ``encompasses trades, limit orders and cancelations'' \citep[arXiv v3, sec.~2.1, p.~4]{catalogue-0503}. The \emph{trade imbalance} is $\mathrm{TI}_k=\sum(b_n-s_n)$, buyer- minus seller-initiated volume (sec.~4.1, p.~16).

\qffield{Conditions} \emph{Evidence.} On one month (April 2010) of consolidated NYSE TAQ quotes and trades for 50 S\&P~500 stocks drawn at random, with 10-second intervals and regressions per half-hour, $\Delta P_k=\beta\,\mathrm{OFI}_k+\epsilon_k$ is linear with an average $R^2$ of 65 percent, against 32 percent for $\mathrm{TI}_k$; with both, the trade imbalance is significant in only 31 percent of subsamples (pp.~8 and~16), and the slope is inversely proportional to depth (abstract, p.~1). Timestamps were rounded to the second and quotes aggregated into an NBBO; Level~1 data only (p.~8). \emph{Names.} Bouchaud et al.\ call signed market-order volume $\sum\varepsilon_n\upsilon_n$ the ``order-flow imbalance'' \citep[chap.~11, sec.~11.4, p.~222]{catalogue-0528}; that is Cont et al.'s trade imbalance. A paper's ``OFI'' must be read from its formula.

\qffield{Practice} Checked on a toy sequence (\nolinkurl{order-flow-checks.py}): a market sell of 200 that is replenished, a cancellation of 300 at the ask and a market buy of 100 that empties the ask give $\mathrm{OFI}=+400$ and $\mathrm{TI}=-100$; the two can differ in sign. The legacy list of imbalance features (trade imbalance, best-quote imbalance as in \hyperref[def:queue-imbalance]{queue imbalance}, multi-level depth imbalance, event imbalance by queue, and pressure aggregated across venues) is a menu of constructions, not results; multi-level and cross-asset OFI are the subject of Cont, Cucuringu and Zhang (a required reading of \texttt{lob-deep-learning}, not read here). Normalising by volume, average daily volume or depth before comparing stocks is interpretation; Cont et al.'s depth result suggests depth.

\qffield{Pitfalls} The legacy notes defined ``order flow imbalance'' in words as net buyer- minus seller-initiated activity, then gave signed trade volume as its formula: that is trade imbalance, not the event-based OFI that the pathway's readings use. Their caveat that imbalance ``can predict the next quote move but fail over longer horizons'' was not found in these sources: Cont et al.\ report contemporaneous fits, not forecasts, and on a subsample of stocks the results held for intervals from 10 quote updates to 10 minutes, the fit improving with the interval (p.~8). Any OFI feature used to forecast must be computed only from events before the forecast time.

\qffield{See also} \hyperref[meth:trade-sign-classification]{trade-sign classification}; \hyperref[def:queue-imbalance]{queue imbalance}; \hyperref[meth:information-driven-bars]{information-driven bars} (imbalance bars, a different use); \hyperref[mod:kyle-single-auction]{Kyle} (impact linear in net order flow); \hyperref[meth:order-flow-impact-regression]{order-flow impact regression}.

\qffield{Further reading} Kolm, Turiel and Westray, and Cont, Cucuringu and Zhang, required readings of \texttt{lob-deep-learning}, are not read here. The legacy file is kept verbatim in \nolinkurl{notes/artefacts/lob-deep-learning/}.

\qffield{Index} order-flow imbalance; OFI; event-based imbalance; trade imbalance; signed volume; best bid and offer; Level 1 data; TAQ; price impact of order book events; depth; imbalance features; multi-level OFI

\qffield{Status} Verification: partly checked 2026-10-16 (against the text layer of \texttt{catalogue-0503} in its arXiv version 3 of 13 April 2011, pp.~1, 4, 8 and~16, downloaded 2026-10-16; the journal version was not checked; and \texttt{catalogue-0528}, Cambridge 2018, p.~222; the toy sequence computed). Stability: draft. Currency: evidence from April 2010 U.S. data.
\end{qfdefinition}
